\section{Memory metadata granularity (Arm F)}
\label{sec:memory}

The conflict unit can also be chosen by relational layout. We compare native
Dolt 2.1.8 JSON metadata, a whole-cell TEXT ablation, and typed rows with fields
\emph{memory, key, slot, value, datatype} and \emph{set-in-version}, whose primary
key is $(memory,key,slot)$. Single-valued status and validity-until use an empty slot;
evidence, applicability and tags use the value as their slot. Competency questions
declare the scalar-versus-set policy and the datatypes. Memory bodies remain
unchanged whole units. The JSON arm uses the initial issues-table migration at
Beads commit \texttt{e1ceedf}; unrelated columns and timestamps are fixed.

Ten explicit divergence templates have twenty instances each, over seeds
1, 7, 42, 99 and 2026. The seeds change values/order but preserve balanced case
counts. Each run uses a fresh, loopback-only real SQL-server fixture, a common
base and two committed branches. Actual \texttt{DOLT\_MERGE} conflict tables are
scored against declared expected outcomes. These are controlled synthetic
experiments, not estimates of the frequency of Memory conflicts in practice.

% Generated by just bench memory-report; do not edit counts.
\begin{table}[ht]\centering\small
\begin{tabular}{lrrrrrr}\toprule
Arm & Units & False units & FP slots & FN slots & Lost & Extra \\ \midrule
Dolt JSON & 100 & 40 & 60 & 0 & 100 & 20 \\
Whole-cell TEXT & 140 & 80 & 140 & 0 & 180 & 60 \\
Typed rows + provenance & 100 & 40 & 40 & 0 & 40 & 20 \\
\bottomrule\end{tabular}
\caption{Arm F: 200 synthetic Memories per seed, 60 true conflict slots. All arms detect the 60 true slots. Units are native conflict rows; false units contain no oracle conflict; FP/FN use semantic slots. Lost/extra counts use automatic output with conflicted units held at base.}
\label{tab:memorymerge}\end{table}

\begin{table}[ht]\centering\small
\begin{tabular}{lrr}\toprule
Competency query & JSON (ms) & Rows (ms) \\ \midrule
Status & 7.825 & 1.227 \\
Tag & 7.627 & 0.789 \\
Traversal & 2.358 & 22.853 \\
\bottomrule\end{tabular}
\caption{Arm F warm SQL latency, pooled median of 35 executions on 2,000 Memories, 10,000 metadata facts and 2,000 Links per seed. One local fixture; these are workload-specific timings.}
\label{tab:memoryquery}\end{table}


Native Dolt JSON already merges edits to independent object keys; the TEXT arm
must therefore not stand in for current Dolt JSON behavior. The native JSON arm
still conflicts on concurrent array changes that the declared set oracle regards
as harmless. Typed rows merge distinct set additions, but storing provenance
inside each row creates false conflicts when two writers assert the same value
with different \texttt{set\_in\_version} values. Consequently the native decision
count and the false native decision count do not improve over JSON in this
balanced trial, although the number of collateral false semantic-slot questions
falls. Provenance reconciliation remains a design
choice, not something the primary key alone settles.

TP/FP/FN counts in Table~\ref{tab:memorymerge} use declared semantic slots;
a conflicted JSON/TEXT Memory expands to all changed metadata keys. This
expansion counts collateral unresolved changes, not separate native prompts.
For loss/extra accounting, every conflicted unit is held at base. Those counts
include deferred valid changes, not irreversible data deletion. Instrument
controls force genuine and false conflicts, missed conflicts, lost/fabricated
facts and invalid writes. Five invalid-type candidates per seed are refused by
the prototype writer before row mutation; native JSON accepts them. The guard
is application validation, not native SHACL enforcement, and a JSON writer could
apply it too.

The status, tag and metadata-to-Link competency queries return independently
constructed identical answers on both layouts. Coverage is checked for every
seed; each timed execution is checked, not merely a count sample. Table~\ref{tab:memoryquery}
reports warm local SQL latency, including result transfer. Indexed rows help
the scalar and membership queries but are slower for the tested traversal join.
This negative result is retained with the full SQL and \texttt{EXPLAIN} plans.
One machine, fixed query order, a synthetic corpus and one indexed layout do not
support general claims about relational or JSON performance.

The artifact is \texttt{benchmark/memory/}. Run \texttt{just bench memory}
with an explicit scratch directory; \texttt{just bench memory-report} generates
the tables from saved measurements. Array order is ignored by the declared set
oracle; ordered-array applications need a different policy. This arm does not
repair the earlier $0/93$ alias blindness or change Arm D's negative result.
